% =====================================================================
% VÍ DỤ CÚ PHÁP - file riêng, KHÔNG nằm trong báo cáo.
% Dùng để xem cách viết công thức, bảng, hình, trích dẫn rồi chép sang các chương.
% Trên Overleaf: Menu -> Main document -> chọn file này để biên dịch thử,
% xem xong thì chọn lại main.tex.
% Nội dung dưới đây chỉ là chữ minh họa, không liên quan bài báo.
%
% LƯU Ý: lớp uetgraduation viết hình và bảng KHÁC cách thông thường:
% chú thích đặt ngay sau \begin{figure} hoặc \begin{table}, không dùng \caption.
% =====================================================================
\documentclass{uetgraduation}
\usepackage{booktabs}
\graphicspath{{hinh/}}
\numberwithin{equation}{chapter}
\mathfont[operator]{Latin Modern Roman}
\newmathbold{\mathbf}{Latin Modern Roman}
\newmathrm{\mathrm}{Latin Modern Roman}

\begin{document}
\setcounter{chapter}{2}
\chapter{Ví dụ cú pháp}

\section{Công thức có đánh số}
Công thức tự đánh số ở lề phải theo chương, như yêu cầu của rubric:
\begin{equation}
	\hat{y}_{t+h} = f(y_1, y_2, \dots, y_t)
	\label{eq:vidu}
\end{equation}
Nhắc lại công thức bằng \verb|\eqref{eq:vidu}| sẽ ra \eqref{eq:vidu}.
Công thức nằm trong dòng viết giữa hai dấu đô la, ví dụ $y_t$ với $t = 1, \dots, T$.

\section{Bảng}
Chú thích bảng tự hiện ở trên đầu bảng (quy định mục 4).
Nhắc bảng bằng \verb|\ref{bang:vidu}|: xem Bảng~\ref{bang:vidu}.

\begin{table}[H]{Ví dụ bảng kết quả (số liệu minh họa)}
	\label{bang:vidu}
	\centering
	\begin{tabular}{lcc}
		\toprule
		\textbf{Mô hình} & \textbf{Bài báo} & \textbf{Tái hiện} \\
		\midrule
		Mô hình A & 1,00 & 1,02 \\
		Mô hình B & 2,00 & 1,98 \\
		\bottomrule
	\end{tabular}
\end{table}

\section{Hình}
Chú thích hình tự hiện ở dưới hình. File ảnh đặt trong thư mục \texttt{hinh/},
viết tên không cần đuôi. Nhắc hình: xem Hình~\ref{hinh:vidu}.

\begin{figure}[H]{Ví dụ hình (ảnh mẫu có sẵn trong LaTeX)}
	\label{hinh:vidu}
	\centering
	\includegraphics[width=0.6\textwidth]{example-image}
\end{figure}

\section{Trích dẫn tài liệu}
Viết \verb|\cite{zhong2025osp}| sẽ ra số trong ngoặc vuông \cite{zhong2025osp}.

\section{Gạch đầu dòng}
\begin{itemize}
	\item Ý thứ nhất.
	\item Ý thứ hai.
\end{itemize}

\subsection{Tiểu mục}
Tiểu mục được đánh ba số. Nhớ để một dòng trống trước tiêu đề mục tiếp theo.

\begin{thebibliography}{9}
\begin{bibsection}{Tiếng Anh}
\bibitem{zhong2025osp}
	Y. Zhong, Y. Ren, G. Cao, F. Li, H. Qi,
	``Optimal starting point for time series forecasting'',
	\textit{Expert Systems with Applications}, 2025.
\end{bibsection}
\end{thebibliography}
\end{document}
